\section{Introduction}
Lexical URL classifiers are attractive for phishing defence because they need only the URL string: no page fetch, DNS lookup or third-party reputation service, which keeps them cheap and usable before a user clicks. Such classifiers routinely report near-ceiling accuracy when trained and tested on random splits of one dataset~\cite{sahingoz2019}. That number answers a narrow question. A deployed detector meets URLs collected by a different process, at a different time and from different hosting ecosystems, and recent work shows that lexical models trained on one public corpus lose much of their accuracy on another~\cite{rashid2024,kanon2026,misiek2026}. In a security setting the failure mode matters: a model that transfers poorly either misses phishing or floods analysts with false alarms on benign sites it has never seen.

Two remedies are common. Domain adaptation uses (unlabelled) target data to align the source and target distributions~\cite{rashid2024}. Cross-dataset feature selection keeps only features that are important on several datasets~\cite{kanon2026,misiek2026}. Both consult the target distribution before evaluation, so a reported gain partly reflects access to the very population on which the model is scored. In addition, URL-level random splits let URLs of the same registered domain fall on both sides of the train/test boundary, which rewards memorising domains rather than learning transferable cues.

This paper studies the stricter \emph{source-only} question: when nothing about the target is available, does choosing features by the \emph{stability} of their importance across held-out domains of the source improve transfer, and does a smaller feature set reduce cost? We make three contributions, which we frame as a carefully controlled empirical study rather than a new algorithm.
\begin{enumerate}
\item A curated, \emph{domain-disjoint} benchmark built from two public sources, EBBU-2017~\cite{sahingoz2019} and HISPAR--Phishstats (HP)~\cite{rashid2024}: registered domains shared by the two sources are removed symmetrically, and every split is group-disjoint by registered domain.
\item A transparent comparison of two source-only selectors built on the same permutation-importance (PI) evidence---ranking by mean PI (MEAN-PI) or by cross-fold stability (STABLE-PI)---at matched feature budgets, for random forest (RF)~\cite{breiman2001} and LightGBM~\cite{ke2017}, in both transfer directions, with nested selection, thresholds locked on source validation data and domain-cluster bootstrap intervals.
\item A measured account of what feature reduction buys at inference time (end-to-end latency, separated into parsing, feature extraction and prediction) and of how the models' SHAP attributions shift between source and target.
\end{enumerate}
The protocol, code and all analysis decisions were frozen and hashed before any test label was read; null and negative findings are reported as obtained.
